% TOP_PROBLEM: {"id": 2303033, "title": "Research Problems in Function Theory — Problem 3.33", "category_id": 8, "category": {"id": 8, "name": "analysis", "display_name": "Analysis", "description": "Limits, continuity, calculus, and function theory.", "slug": "analysis", "order_index": 8, "created_at": "2026-07-31T15:26:25.671Z"}, "status": "open", "problem_number": "AMR-022-3033", "set_id": 13}
% FIRST_POSTED: 2026-10-02
% ENTRY_KIND: statement_only
% UnsolvedMath ID 2303033; source catalogue code AMR-022-3033
% !TeX program = lualatex
% Definition and sources only; this file is not an LLM solution attempt.
\documentclass[11pt,a4paper]{article}
\usepackage[margin=25mm]{geometry}
\usepackage{fontspec}
\setmainfont{lmroman10-regular.otf}[BoldFont=lmroman10-bold.otf,
 ItalicFont=lmroman10-italic.otf,BoldItalicFont=lmroman10-bolditalic.otf]
\usepackage{amsmath,amssymb,mathtools}
\usepackage{unicode-math}
\setmathfont{latinmodern-math.otf}
\AtBeginDocument{\let\setminus\smallsetminus}
\usepackage{xurl,hyperref}
\hypersetup{colorlinks=true,urlcolor=blue}
\setcounter{secnumdepth}{0}
\setlength{\parindent}{0pt}
\setlength{\parskip}{5pt}
\setlength{\emergencystretch}{3em}
\begin{document}
\section*{Research Problems in Function Theory — Problem 3.33}\label{2303033}
{\small UnsolvedMath ID 2303033 · AMR-022-3033 · Analysis · AMR Open Problem Lists}
\subsection{Definitions and mathematical statement}
Let $\Omega\subset\mathbb R^n$ be bounded and open, $n\ge2$, and $y\in\partial\Omega$. For boundary data $f$, the Perron--Wiener--Brelot method defines upper and lower harmonic envelopes using superharmonic majorants and subharmonic minorants with the required boundary limits. Call $f$ resolutive when these envelopes coincide in a finite harmonic function, denoted $H_f^\Omega$.
The point $y$ is $B$-regular for $\Omega$ if, whenever resolutive real boundary data $f$ are bounded in some neighbourhood of $y$ on $\partial\Omega$, their solution $H_f^\Omega$ is bounded in some neighbourhood of $y$ inside $\Omega$. Explicitly, for each such $f$, an $R>0$ with
$\sup_{\partial\Omega\cap B(y,R)}|f|<\infty$ must imply the existence of $r>0$ with
$\sup_{\Omega\cap B(y,r)}|H_f^\Omega|<\infty$.
Call $y$ locally $B$-regular if there is a sequence $r_m\downarrow0$ such that $y$ is $B$-regular for every truncated open set $\Omega\cap B(y,r_m)$, with its own complete boundary and resolutive data.
Local $B$-regularity implies $B$-regularity. Is the converse true for every $\Omega$ and $y$? The neighbourhood in the definition may depend on $f$; no uniform bound over all boundary data is required.
\subsection{Short English statement}
Does boundedness regularity for the Dirichlet problem at a boundary point persist on arbitrarily small ball truncations of the domain?
\subsection{Sources}
\begin{itemize}
\item[S1] ULAM.AI and the UnsolvedMath Contributors, supplied problems.json, record 2303033 (AMR-022-3033), AMR Open Problem Lists. Catalogue identity and source transcription; CC BY 4.0.
\url{https://huggingface.co/datasets/ulamai/UnsolvedMath}.
\item[S2] Hayman - Research Problems in Function Theory (2018), Problem 3.33. Original source indexed by AMR-022-3033.
\url{https://arxiv.org/abs/1809.07200}.
\item[S3] A. Sadi, Some types of regularity for the Dirichlet problem, Section 6. Primary definition and converse question.
\url{https://www.cambridge.org/core/services/aop-cambridge-core/content/view/620E7F72E60E8BB1A78F41A02E75F9F5/S0027763000004013a.pdf/some_types_of_regularity_for_the_dirichlet_problem.pdf}.
\end{itemize}
\end{document}
